\documentclass[10pt,a4paper]{amsart}
\usepackage[margin=19mm]{geometry}
\usepackage{amsmath,amssymb,mathtools}
\usepackage[scr=rsfs,cal=euler]{mathalfa}
\usepackage{xfrac}
\usepackage[unicode,hidelinks,bookmarks=false]{hyperref}
\newcommand{\ie}{i.e.\ }
\newcommand{\cf}{cf.\ }
\newcommand{\resp}{respectively\ }
\newcommand{\Subsection}{\subsection}
\providecommand{\nobreakdash}{\nobreak\hbox{-}}
\setlength{\emergencystretch}{2em}
\multlinegap0pt
\usepackage{amssymb}
\usepackage{float}
\newtheorem{lemma}{Lemma}
\newtheorem{proposition}{Proposition}
\newcommand{\bbp}{{\mathbb P}}
\newcommand{\bbr}{{\mathbb R}}
\newcommand{\bbrn}{{\bbr^n}}
\newcommand{\cQ}{{\mathcal Q}}
\title{Asymptotic behavior of small solutions to the Vlasov--Klein--Gordon system in high dimensions}
\author{Ho Lee}
\date{Source: arXiv:2603.28438v2 (CC BY 4.0)}
\begin{document}
\maketitle

\noindent\emph{Source and licence.} Asymptotic behavior of small solutions to the Vlasov--Klein--Gordon system in high dimensions. Version arXiv:2603.28438v2 (CC BY 4.0), distributed by arXiv under the Creative Commons Attribution 4.0 International licence, http://creativecommons.org/licenses/by/4.0/, as recorded in the arXiv licence metadata for this exact version and shown on the abstract page https://arxiv.org/abs/2603.28438v2. Full licence notice: \texttt{licenses/arxiv-2603.28438-CC-BY-4.0.txt}.\par\smallskip

\noindent\textit{Editorial scope.} Counted unit: the article's lemma homo (source0.tex lines 1113--1119). The article's complete original proof of it is delivered with it (source0.tex lines 1120--1188, one \texttt{proof} environment of 69 lines). No mathematics is written, repaired or re-derived: the statement and the proof are the article's own lines, in the article's own notation, and no retained line is substituted or re-flowed.  Retained uncounted as the source-local context the proof needs: lines 211--213; lines 398--403; lines 539--546; lines 608--630; lines 706--712; lines 796--799; lines 1097--1099.  Nothing was omitted from any retained span, so the package carries no omission record.  No retained line contains a citation, so the wrapper carries no bibliography and no bibliography entry.  No retained line contains a figure environment, an image-inclusion command or drawing code, so no image is delivered and the package declares no asset dependency.  Editorial formatting only, in addition: an AMS \texttt{amsart} preamble that restates the article's own declarations and macros used by the retained lines, namely the environments \texttt{lemma}, \texttt{proposition} on separate counters, together with the standard packages those lines need.  The retained environments print in this document as Lemma 1, Proposition 1 in order of appearance, whereas the article prints the same items with the numbering of its own full document.  The complete native source of the article as distributed in the arXiv e-print for this exact version is bundled unchanged as \texttt{sources/197438/source0.tex}.
\begin{align}
\tau = \sqrt{ t^2 - | x |^2 } , \qquad y = x . \label{tau}
\end{align}

\begin{lemma} \label{lem intchif}
Let $ f $ and $ h $ be regular functions satisfying $ T_\phi f = h $ defined in $ \bigcup_{ 1 \leq \tau < T } H_\tau \times \bbr^n_v $. Then we have
\begin{align*}
& \int_{ H_\tau } \hat{ e } ( | f |) \, d \mu_{ H_\tau } \leq \int_{ H_1 } \hat{ e } ( | f | ) \, d \mu_{ H_1 } + \int_1^\tau \int_{ H_\lambda } \int_\bbrn \left( | h | + | \nabla_x \phi | | f | \right) d v \, d \mu_{ H_\lambda } d \lambda .
\end{align*}
\end{lemma}

\begin{proposition} \label{prop KSphi}
Let $ \phi $ be a regular function defined in $ \bigcup_{ 1 \leq \tau < T } H_\tau $ for some $ T \in ( 1 , \infty ] $. Then we have
\begin{align}
| \phi | ( t , x ) & \lesssim \frac{ 1 }{ t^\frac{ n }{ 2 } } E_{ \lfloor n / 2 \rfloor +1 }^{ \frac12 } ( \phi ) ( \tau ) , \label{KSphi} \\
| \partial \phi | ( t , x ) & \lesssim \frac{ 1 }{ t^{ \frac{ n }{ 2 } - 1 } \tau } E_{ \lfloor n / 2 \rfloor +1 }^{ \frac12 } ( \phi ) ( \tau ) , \label{KSphiimprove}
\end{align}
for $ t $, $ x $ and $ \tau $ satisfying \eqref{tau}.
\end{proposition}

\begin{lemma} \label{lem basic}
The following commutation rules hold:
\begin{enumerate}
\item For any $ \hat{ Z }_a \in \hat{ \bbp } $, the free transport operator $ T $ satisfies
\begin{align*}
\left[ T , \hat{ Z }_a \right] = 0 .
\end{align*}
\item For any $ \hat{ Z }_a , \hat{ Z }_b \in \hat{ \bbp } $, there exist constant coefficients $ C_{ a b }^c $ such that
\begin{align*}
\left[ \hat{ Z }_a , \hat{ Z }_b \right] = \sum_{ \hat{ Z }_c \in \hat{ \bbp } } C_{ a b }^c \hat{ Z }_c .
\end{align*}
\item Concerning the derivatives $ \partial_{ v^i } $ we have
\begin{align*}
& \left[ \hat{ \partial }_t , \partial_{ v^i } \right] = 0 , \\
& \left[ \hat{ \partial }_{ x^i } , \partial_{ v^j } \right] = 0 ,\\
& \left[ \hat{ \Omega }_{ 0 i } , \partial_{ v^j } \right] = - \frac{ v^j }{ v^0 } \partial_{ v^i } , \\
& \left[ \hat{ \Omega }_{ i j } , \partial_{ v^k } \right] = - \delta^i_k \partial_{ v^j } + \delta^j_k \partial_{ v^i } , \\
& \left[ \hat{ \Omega }_{ 0 i } , v^j \partial_{ v^j } \right] = \frac{ 1 }{ v^0 } \partial_{ v^i } , \\
& \left[ \hat{ \Omega }_{ i j } , v^k \partial_{ v^k } \right] = 0 ,
\end{align*}
for any $ i , j , k = 1 , \dots , n $.
\end{enumerate}
\end{lemma}

\begin{lemma} \label{lem TphiZhat3}
For any $ \hat{ Z }_A \in \hat{ \bbp }^{ | A | } $ with $ | A | \geq 1 $, we have
\begin{align*}
\left[ T_\phi , \hat{ Z }_A \right] = \sum_{ \substack{ | B | + | C | \leq | A | + 1 \\ 1 \leq | C | \leq | A | } } \sum_{ 0 \leq \mu \leq n } Q^{ \mu B C }_{ A } \partial_{ x^\mu } ( Z_B \phi ) \hat{ Z }_C ,
\end{align*}
for some $ Q^{ \mu B C }_{ A } \in \cQ $.
\end{lemma}

\begin{align}
E_N ( \phi ) ( \tau ) & < 2 \varepsilon , \label{bootsphi} \\
\hat{ E }_{ N , 1 } ( f ) ( \tau ) & < 2 \varepsilon \tau^\delta , \label{bootsf}
\end{align}

\begin{align} \label{TphiH}
T_\phi F_1 + R_1 F_1 = 0 ,
\end{align}

\begin{lemma} \label{lem homo}
The homogeneous part $ F_1 $ satisfies
\begin{align}
\hat{ E }_n ( F_1 ) ( \tau ) \lesssim \varepsilon \tau^\delta , \label{homo}
\end{align}
where $ \delta = 0 $ for $ n \geq 5 $, and $ \delta = \varepsilon^{ 1 / 4 } $ for $ n = 4 $.
\end{lemma}

\begin{proof}
Applying $ \hat{ Z }_D $ to the equation \eqref{TphiH} we obtain
\begin{align}
T_\phi \hat{ Z }_D F_1 = - \hat{ Z }_D \left( R_1 F_1 \right) + \left[ T_\phi , \hat{ Z }_D \right] F_1 . \label{TphiZF}
\end{align}
Concerning the first term on the right hand side we notice that $ R_1 $ is given by
\begin{align*}
R_1 = \sum_{ | B | + | C | \leq | A | + 1 } \sum_{ 0 \leq \mu \leq n } Q^{ \mu B C }_{ A } \partial_{ x^\mu } ( Z_B \phi ) ,
\end{align*}
with multi-indices $ A $ and $ C $ satisfying
\begin{align*}
\left\lfloor \frac{ N }{ 2 } \right\rfloor - n +1 \leq | A | \leq N , \qquad \left\lfloor \frac{ N }{ 2 } \right\rfloor - n +1 \leq | C | \leq | A | .
\end{align*}
Then, we obtain for any $ | D | \leq n $,
\begin{align*}
| \partial_{ x^\mu } ( Z_{ D } Z_B \phi ) | & \lesssim \frac{ 1 }{ t^{ \frac{ n }{ 2 } - 1 } \tau } E_{ \lfloor n / 2 \rfloor + 1 }^{ \frac12 } ( Z_{ D } Z_B \phi ) \\
& \lesssim \frac{ \varepsilon^{ \frac{ 1 }{ 2 } } }{ t^{ \frac{ n }{ 2 } - 1 } \tau } ,
\end{align*}
by Proposition \ref{prop KSphi} and the bootstrap assumption \eqref{bootsphi}, since $ N \geq 5 n + 2 $ and
\begin{align*}
\left\lfloor \frac{ n }{ 2 } \right\rfloor + 1 + | D | + | B | & \leq \left\lfloor \frac{ n }{ 2 } \right\rfloor + 1 + n + N + 1 - \left\lfloor \frac{ N }{ 2 } \right\rfloor + n - 1 \\
& \leq \left\lfloor \frac{ 5 n + 2 }{ 2 } \right\rfloor + N - \left\lfloor \frac{ N }{ 2 } \right\rfloor \\
& \leq N .
\end{align*}
We recall that $ \hat{ Z }_{ D } Q^{ \mu B C }_{ A } \in \cQ $, so that $ | \hat{ Z }_{ D } Q^{ \mu B C }_{ A } | \lesssim t $, and note that $ Z_{ D } $ commutes with $ \partial_{ x^\mu } $ as in Lemma \ref{lem basic}. Now, we can conclude that for any $ | D | \leq n $,
\begin{align*}
| \hat{ Z }_D ( R_1 F_1 ) | & \lesssim \frac{ \varepsilon^{ \frac{ 1 }{ 2 } } }{ t^{ \frac{ n }{ 2 } - 2 } \tau } \sum_{ | D_1 | \leq n } | \hat{ Z }_{ D_1 } F_1 | .
\end{align*}
For the second term on the right hand side of \eqref{TphiZF} we use Lemma \ref{lem TphiZhat3} to have
\begin{align*}
\left[ T_\phi , \hat{ Z }_D \right] = \sum_{ \substack{ | B | + | C | \leq | D | + 1 \\ 1 \leq | C | \leq | D | } } \sum_{ 0 \leq \mu \leq n } Q^{ \mu B C }_{ D } \partial_{ x^\mu } ( Z_B \phi ) \hat{ Z }_C ,
\end{align*}
for some $ Q^{ \mu B C }_{ D } \in \cQ $. Since $ | D | \leq n $, we obtain by the same arguments as above
\begin{align*}
\left| \left[ T_\phi , \hat{ Z }_D \right] F_1 \right| & \lesssim \frac{ \varepsilon^{ \frac{ 1 }{ 2 } } }{ t^{ \frac{ n }{ 2 } - 2 } \tau } \sum_{ | D_1 | \leq n } | \hat{ Z }_{ D_1 } F_1 | .
\end{align*}
Now, we apply Lemma \ref{lem intchif} to obtain
\begin{align*}
& \int_{ H_\tau } \hat{ e } ( | \hat{ Z }_D F_1 | ) \, d \mu_{ H_\tau } - \int_{ H_1 } \hat{ e } ( | \hat{ Z }_D F_1 | ) \, d \mu_{ H_1 } \\
& \lesssim \int_1^\tau \int_{ H_\lambda } \int_\bbrn | \hat{ Z }_D ( R_1 F_1 ) | + \left| \left[ T_\phi , \hat{ Z }_D \right] F_1 \right| + | \nabla_x \phi | | \hat{ Z }_D F_1 | \, d v \, d \mu_{ H_\lambda } d \lambda \\
& \lesssim \int_1^\tau \int_{ H_\lambda } \frac{ \varepsilon^{ \frac{ 1 }{ 2 } } }{ t^{ \frac{ n }{ 2 } - 2 } \lambda } \sum_{ | D_1 | \leq n } \hat{ e } ( | \hat{ Z }_{ D_1 } F_1 | ) \, d \mu_{ H_\lambda } d \lambda \\
& \lesssim \varepsilon^{ \frac{ 1 }{ 2 } } \int_1^\tau \frac{ 1 }{ \lambda^{ \frac{ n }{ 2 } - 1 } } \hat{ E }_n ( F_1 ) \, d \lambda ,
\end{align*}
which implies that
\begin{align}
\hat{ E }_n ( F_1 ) ( \tau ) \leq \hat{ E }_n ( F_1 ) ( 1 ) + C \varepsilon^{ \frac{ 1 }{ 2 } } \int_1^\tau \frac{ 1 }{ \lambda^{ \frac{ n }{ 2 } - 1 } } \hat{ E }_n ( F_1 ) ( \lambda ) \, d \lambda . \label{FGronwall}
\end{align}
Since $ F_1 ( 1 ) = F ( 1 ) $ and $ \hat{ E }_n ( F ) ( 1 ) \lesssim \hat{ E }_{ N + n } ( f ) ( 1 ) < \varepsilon $, we can find $ C_{ 1 } > 0 $ and $ \tau_{ 1 } > 1 $ such that
\begin{align}
\hat{ E }_n ( F_1 ) ( 1 ) \leq C_{ 1 } \varepsilon , \qquad \hat{ E }_n ( F_1 ) ( \tau ) \leq
\left\{
\begin{aligned}
& 2 C_{ 1 } \varepsilon & ( n \geq 5 ) , \\
& 2 C_{ 1 } \varepsilon \tau^\delta & ( n = 4 ) ,
\end{aligned}
\right. \label{Flocal}
\end{align}
for $ 1 \leq \tau < \tau_{ 1 } $. Then, the inequality \eqref{FGronwall} implies that
\begin{align*}
\hat{ E }_n ( F_1 ) ( \tau ) \leq C_{ 1 } \varepsilon + C
\left\{
\begin{aligned}
& C_{ 1 } \varepsilon^{ \frac{ 3 }{ 2 } } & ( n \geq 5 ) , \\
& C_{ 1 } \varepsilon^{ \frac{ 3 }{ 2 } } \delta^{ - 1 } \tau^\delta & ( n = 4 ) ,
\end{aligned}
\right.
\end{align*}
for $ 1 \leq \tau < \tau_{ 1 } $. We take $ \delta = \varepsilon^{ 1 / 4 } $ to conclude that \eqref{Flocal} holds for all $ \tau \geq 1 $.
\end{proof}

\end{document}
